\centerline{\bf \Huge Можно ли? II}

\begin{flushright}
        \scriptsize{Можно все, что хочешь, когда хочешь только то, что следует.}
\end{flushright}

\problem{\textbf{1}} Можно ли в прямоугольной таблице расставить натуральные числа так, чтобы в каждом столбце сумма чисел была больше 100, а в каждой строке – меньше 5?

\problem{\textbf{2}} Можно ли подобрать такие два последовательных натуральных числа, что суммы цифр каждого числа равны?

\problem{\textbf{3}} Можно ли разрезать квадрат на квадратики двух размеров так, чтобы маленьких было столько же, сколько и больших?

\begin{flushright} \textbf{3 балла}\\ 
\end{flushright}

\problem{\textbf{4}} На школьном спектакле все 25 мест в первом ряду заняты
школьниками. Известно:

1) Никакие две девочки в этом ряду не сидят рядом;

2) Рядом с каждым мальчиком сидит ещё хотя бы один мальчик;

3) Всего в первом ряду сидят 9 девочек.

Могло ли так оказаться, что на центральном месте в ряду сидит мальчик?

\problem{\textbf{5}} Имеется 12 сосисок длиной 13 см каждая. Их
требуется разделить между 13 котятами, 13 кошками и 13 котами так, чтобы каждому
котенку достался кусок сосиски длиной 3 см, каждой кошке — кусок сосиски длиной
4 см, а каждому коту кусок сосиски длиной 5 см. Можно ли это сделать?
\begin{flushright} \textbf{4 балла}\\ 
\end{flushright}


\problem{\textbf{6}} Начальник выписал на доску в ряд некоторые целые числа от 0 до
999. В итоге получилось длинное число. НеНачальник записал на свою часть доски все оставшиеся
целые числа из этого же диапазона, в итоге получилось второе длинное число. Могли ли эти
два длинных числа совпасть?


\problem{\textbf{7}} На территории страны, имеющей форму квадрата со стороной 1000 км, находится 51 город. Страна располагает средствами для прокладки 11000 км. Сможет ли правительство страны соединить сетью дорог все свои города?

\begin{flushright} \textbf{5 баллов}\\ 
\end{flushright}